\chapter{Conclusion}
\label{ch:conclusion}

% Conclude by answering the question posed in the introduction. Separate
% established results, limitations, and future work; do not introduce a new
% theorem or experiment here.

We defined reachability as the reflexive transitive closure in
\cref{eq:transitive-closure}, proved its transitivity in
\cref{thm:transitivity}, and computed it with \cref{alg:warshall}. The graph and
matrix in \cref{fig:example-graph,tab:reachability} provided a transparent check
of the definitions and algorithm.

The example deliberately leaves many questions open, including path recovery,
sparse-graph algorithms, and incremental updates. Its main purpose is structural:
it demonstrates how a thesis can connect a research question to definitions,
claims, evidence, limitations, and correctly cited sources without adopting a
publisher-style layout.
